\scanpage{140}
\(S\) schéma normal intègre, pt générique \(\eta\), \(k(\eta)=K\).

\uline{Isomotif} (\(\frac12\) simple) \uline{lisse} sur \(S\) :
isomotif semi-simple \inserted{\(M\)} sur \(S\), tel que pour tout
\(l\in P\), \(l\ne\operatorname{car}K\), \(T_l(M)\) soit non ramifié rel. à
\(S_l=S[1/l]\).
Conjecturalement, il suffit que ce soit vrai pour \uline{un} \(l\)
premier à toutes les car. résiduelles
(\(\Longrightarrow\) il suffit que ce soit vrai pour \uline{deux} \(l\ne l'\)).

Motif (\unsure{isosemisimple|iso semisimple}) lisse sur \(S\)
(si \(\operatorname{car}K=0\)) :
\[
=\{M,\{E_l\},\{\alpha_l\}\}
\]
où

a) \(M\) un isomotif lisse sur \(S\)

b) \(\forall l\), \(E_l\) un faisceau \(l\)-adique \unsure{ct tordu|ctordu}
sur \(S_l=S[1/l]\) (ou donc sur \(K\), mais non ramifié rel. à \(S[1/l]\))

c) \(\forall l\), \(\alpha_l:E_l\otimes_{Z_l}Q_l\xrightarrow{\sim}T_l(M)\)

avec un axiome : ``pour presque tout \(l\), \(\unsure{u_l|\alpha_l}\)
provient d'un isomorphisme sur \(Z_l\) --- pas seulement à torsion près''.

Ceci implique que pour presque tout \(l\), \(E_l\) est sans torsion.

Homomorphisme de motifs lisses sur \(S\)
\[
(M,\{E_l\},\{\alpha_l\})\unsure{\mathbin{\text{---}}|\longrightarrow}(\unsure{M|M'},\{E'_l\},\{\alpha'_l\})
\]
\inserted{\(\unsure{u\to M}\)}
donné par \((u,\{u_l\})\), avec \(u:M\mathbin{\text{---}}M'\) un hom de
isomotifs, \unsure{à} et \(u_l:E_l\mathbin{\text{---}}\unsure{E_l|E'_l}\) un hom de
faisceaux \(l\)-adiques tels que pour \(\forall l\), on ait commutativité du
\diagram{\[
\begin{tikzcd}
E_l\otimes_{Z_l}Q_l\arrow[r,no head,"\alpha_l"]\arrow[d,no head,"u_l\otimes_{Z_l}Q_l"']&
T_l(M)\arrow[d,"\unsure{T_l(u)|T_l(M)}"]\\
E'_l\otimes_{Z_l}Q_l\arrow[r,no head,"\alpha'_l"']&T_l(M')
\end{tikzcd}
\]}

\uline{N.B.}\quad Le foncteur \deleted{\unsure{Hom}}
\[
\text{Motifs sur }K\longleftarrow\text{Motifs lisses sur }S
\]
est pleinement fidèle. Les Hom sont des modules de type fini sur \(Z\).
La catégorie des motifs lisses sur \(S\) est abélienne
(on prend \uline{noyaux}, conoyaux séparément pour les composantes \(M,E_l\)).
\deleted{Le foncteur} \inserted{L'image de}
\deleted{\unsure{Mot}}
(Motifs lisses sur \(S\)) dans (Motifs sur \(K\))
est une ss-catégorie abélienne \uline{épaisse}.

Le foncteur (Motifs\inserted{ lisses})
\(\longrightarrow\) (Isomotifs\inserted{ lisses}) est exact et surjectif
% archive foliation: 139
